\scanpage{89}
% archive foliation: 88
\uline{1.4 suite}

Choisissons d'abord un splittage de \inserted{$\mathcal L$ i.e. \unsure{de|un}}
$\mathfrak g$, en
\[
\mathfrak g\simeq\Omega+\underset{\substack{\Vert\\O_X}}{\Omega^0}+\Omega^{-1}
\]
d'où, posant $\omega=\underline\Omega^1_{X/S}$
\[
\mathfrak g\otimes\omega\simeq\Omega\otimes\omega+\omega+\Omega^{-1}\otimes\omega
\]
\deleted{\unsure{On va avoir|On va à une}}
\par \unsure{Choisissons-en|Choisissons} loc. une section $\delta$ de $\Omega$ comme
base. Une connexion (sans courbure)
$c_0$ est obtenue par
\diagram{\begin{tikzcd}\mathfrak g\arrow[r,"c_0"]&\mathfrak g\otimes\omega\end{tikzcd}}
\[
c_0(\delta)=c_0(1)=c_0(\delta^{-1})=0
\]
i.e.
\[
c_0\bigl(\underset{\substack{\downarrow\\(\lambda,\mu,\nu)}}{\underbrace{\lambda\delta+\mu.1+\nu\delta^{-1}}}\bigr)
=\underset{\substack{\downarrow\\(d\lambda,d\mu,d\nu)}}{\delta\,d\lambda+1\,d\mu+\delta^{-1}d\nu}
\]

Une connexion générale s'obtient alors en
ajoutant une section de
\[
\operatorname{D\acute er}(\mathfrak g,\mathfrak g)\otimes\omega\simeq\mathfrak g\otimes\omega
\]
(\uline{NB} toute dérivation de $\mathfrak g$ est intérieure!)

Soit $\inserted{\bar\omega_*:=}\bar\omega_2\oplus\bar\omega_1\oplus\bar\omega_0
\in\Gamma\mathfrak g\otimes\omega$\quad i.e.
\[
\begin{aligned}
\bar\omega_2&\in\Gamma(\Omega\otimes\omega)\\
\bar\omega_1&\in\Gamma(\omega)\\
\bar\omega_0&\in\Gamma(\Omega^{-1}\otimes\omega)
\end{aligned}
\]
on aura donc pour connexion $C^{\deleted{\unsure{\bar\omega_*|\bar\omega_2}}}_{\bar\omega_*}$
définie par $c_0$ et $\bar\omega_*$
\[
C_{\bar\omega_*}\bigl(\underbrace{\lambda\delta+\mu.1+\nu\delta^{-1}}_{(\lambda,\mu,\nu)}\bigr)
\overset{?}{=}\deleted{\unsure{\delta\,d\lambda+1\,d\mu+\delta^{-1}d\nu}}
\]
\[
+
\]

